\section{Projects for Chapter 5: Discrete Random Variables and Probability Laws}

This chapter is ideal for interactive graphics because a discrete law has several equivalent representations: a mechanism, a mapping from outcomes to numbers, a probability mass function, a cumulative distribution function, and numerical summaries such as expectation and variance.

\subsection*{Project: outcome-to-random-variable mapper}

\begin{examplebox}
\textbf{Prompt to give an AI coding assistant}

\begin{quote}\small
Create an interactive application that begins with a finite sample space and lets the user define a random variable by assigning a numerical value to each elementary outcome. Group all outcomes that receive the same value and construct the support, probability mass function, and cumulative distribution function automatically.

When the user clicks a value of the random variable, highlight its complete inverse image in the sample space. Display
\[
P(X=x),\qquad F_X(t)=P(X\leq t),\qquad \E[X],\qquad \Var(X).
\]
Include preset mappings for the number of heads in coin tosses, the sum of two dice, a payoff, and an indicator variable. Verify that PMF values sum to one and that CDF jumps equal point masses.
\end{quote}
\end{examplebox}

\subsection*{Project: discrete-law explorer}

The application should include Bernoulli, binomial, geometric, and Poisson laws. Parameters should be controlled by sliders, while support, PMF, CDF, mean, and variance update immediately.

\begin{examplebox}
\textbf{Prompt to give an AI coding assistant}

\begin{quote}\small
Build a discrete probability-law explorer for Bernoulli, binomial, geometric, and Poisson distributions. For each law, show the random mechanism, parameter restrictions, support, PMF, CDF, expectation, and variance.

Use linked PMF and CDF plots. Clicking a PMF bar should highlight the corresponding jump in the CDF. Add a Monte Carlo panel with controls for one sample and many samples. Overlay empirical relative frequencies on the exact PMF and report the maximum absolute difference.

For the binomial law, allow \(n\) and \(p\) to change and verify that \(n=1\) gives a Bernoulli law. For the geometric law, state clearly whether the support begins at zero or one and keep that convention consistent.
\end{quote}
\end{examplebox}

\subsection*{Possible extension: expectation as a balance point}

Draw the PMF as masses placed on a number line and mark \(\E[X]\) as the balancing location. The student should be able to move probability mass while preserving total mass one and observe how expectation and variance change.

\begin{summarybox}
A good discrete-law application links four views:
\[
\text{mechanism}
\longrightarrow
\text{random variable}
\longrightarrow
\text{PMF and CDF}
\longrightarrow
\text{moments}.
\]
\end{summarybox}
